% Changes since the previous version (#428) — the plain-language
% companion's change log, one entry per published release, newest first,
% the same releases as changes-methods.tex. Written in the sitting that
% closes a release, from a diff of the published tree against the previous
% published copy. Format guard: py/test_changes.py.
\clearpage
\section*{Changes since the previous version}
\addcontentsline{toc}{section}{Changes since the previous version}
What changed in this companion since the copy published before it, newest
first.

\subsection*{1.1 (43) --- 30 September 2026}
Against the first public copy, published on 23 September 2026.
\begin{itemize}
  \item \textbf{A new chapter: ``Floors: how deep a reading can see''.}
  It sits after the Waveform Matrix chapter and explains what a floor is,
  where floors come from, why the largest one decides, what a reading
  near a floor becomes, why averaging lowers noise but not the analysis's
  own residue, and why a reading above every floor can still be something
  other than the pedal. Every earlier chapter's floor paragraph points to
  it.
  \item \textbf{Chord IMD.} When the played notes moved during a capture,
  the record no longer writes off every invented tone. It works out how
  much of the notes' spill reaches each one, reads the tones that stand
  clear of it, and marks only the others unread.
  \item \textbf{Harmonic Distortion.} The very last point of each sweep,
  where the sweep fades out, is left out of every reading, and the
  simulated pedals the checks use are run the way the app's Simulation
  Mode runs them.
  \item \textbf{Transfer Curve.} One more drive note, in the top octave.
  \item \textbf{Waveform Matrix.} Two nearly identical interface settings
  could give the largest grid two different sets of quiet levels; they
  now give the same one, and grids recorded before keep theirs.
  \item \textbf{Every chapter.} The chapters now say that everything in
  this companion describes a pedal that has settled. The sections on how
  you can check that the app does what it says are broken into short
  named paragraphs, several passages are reworded, and spelling is
  American throughout.
\end{itemize}
